% !TEX program = xelatex
% !TEX encoding = UTF-8
% Path, Site, Design --- multi-model code evaluation.
% Build: xelatex shared-compare-design.tex && xelatex shared-compare-design.tex
\documentclass[
  11pt,
  a4paper,
  oneside,
  numbers=noenddot,
  abstract=true,
  headings=normal,
  DIV=11,
  BCOR=0mm,
]{scrartcl}

\usepackage{fontspec}
\usepackage{unicode-math}

\setmainfont{texgyrepagella}[
  Extension=.otf,
  UprightFont=*-regular,
  ItalicFont=*-italic,
  BoldFont=*-bold,
  BoldItalicFont=*-bolditalic,
  Ligatures=TeX,
]
\setsansfont{texgyreheros}[
  Extension=.otf,
  UprightFont=*-regular,
  ItalicFont=*-italic,
  BoldFont=*-bold,
  BoldItalicFont=*-bolditalic,
  Scale=MatchLowercase,
  Ligatures=TeX,
]
\setmonofont{texgyrecursor}[
  Extension=.otf,
  UprightFont=*-regular,
  ItalicFont=*-italic,
  BoldFont=*-bold,
  BoldItalicFont=*-bolditalic,
  Scale=MatchLowercase,
]
\setmathfont{texgyrepagella-math.otf}

\usepackage{setspace}
\setstretch{1.07}
\usepackage{microtype}
\usepackage{parskip}
\KOMAoptions{parskip=half-}
\setkomafont{disposition}{\rmfamily\bfseries}
\setkomafont{section}{\Large}
\setkomafont{subsection}{\large}
\setkomafont{subsubsection}{\normalsize}

\RedeclareSectionCommand[
  beforeskip=1.5em plus 0.25em minus 0.1em,
  afterskip=0.55em plus 0.1em,
]{section}
\RedeclareSectionCommand[
  beforeskip=1.1em plus 0.2em minus 0.1em,
  afterskip=0.4em plus 0.1em,
]{subsection}

\usepackage{amsthm}
\usepackage{booktabs}
\usepackage{tabularx}
\usepackage{array}
\usepackage{graphicx}
\usepackage{xcolor}
\usepackage{listings}
\usepackage{caption}
\usepackage{fancyhdr}
\usepackage{etoolbox}

\definecolor{rulegray}{HTML}{4A4A4A}
\definecolor{linkblue}{HTML}{1A5276}
\definecolor{codecode}{HTML}{1C1C1C}
\definecolor{codebg}{HTML}{F4F4F2}
\definecolor{codeframe}{HTML}{D0D0C8}
\definecolor{accent}{HTML}{2C3E50}
\definecolor{muted}{HTML}{666666}

\captionsetup{font=small, labelfont={bf,sf}, labelsep=period, skip=0.45em}

\lstset{
  basicstyle=\ttfamily\footnotesize\color{codecode},
  frame=single,
  framerule=0.4pt,
  rulecolor=\color{codeframe},
  breaklines=true,
  columns=fullflexible,
  keepspaces=true,
  showstringspaces=false,
  aboveskip=0.65em,
  belowskip=0.65em,
  xleftmargin=0.55em,
  xrightmargin=0.55em,
  backgroundcolor=\color{codebg},
  captionpos=b,
  tabsize=2,
}

\usepackage[
  colorlinks=true,
  linkcolor=accent,
  citecolor=accent,
  urlcolor=linkblue,
  bookmarks=true,
  bookmarksnumbered=true,
  pdftitle={Path, Site, Design: Multi-Model Code Evaluation},
  pdfauthor={Design Proposal},
  pdfsubject={Comparison geometry for ranking peer code patches},
  pdfkeywords={model evaluation, code generation, multi-agent, ranking, diffs},
]{hyperref}
\usepackage{bookmark}

\setlength{\headheight}{22pt}
\pagestyle{fancy}
\fancyhf{}
\fancyhead[L]{\small\color{muted}\textit{Path\,$\rightarrow$\,Site\,$\rightarrow$\,Design}}
\fancyhead[R]{\small\color{muted}Multi-Model Code Evaluation}
\fancyfoot[C]{\small\thepage}
\renewcommand{\headrulewidth}{0.3pt}
\renewcommand{\footrulewidth}{0pt}
\renewcommand{\headrule}{\hbox to\headwidth{\color{rulegray}\leaders\hrule height \headrulewidth\hfill}}
\fancypagestyle{plain}{%
  \fancyhf{}%
  \fancyfoot[C]{\small\thepage}%
  \renewcommand{\headrulewidth}{0pt}%
}

\theoremstyle{definition}
\newtheorem{definition}{Definition}[section]
\newtheorem{construction}[definition]{Construction}
\theoremstyle{plain}
\newtheorem{proposition}[definition]{Proposition}
\newtheorem{lemma}[definition]{Lemma}
\newtheorem{corollary}[definition]{Corollary}
\theoremstyle{remark}
\newtheorem{remark}[definition]{Remark}
\newtheorem{example}[definition]{Example}

\newcommand{\PathSiteDesign}{\textbf{path\,$\rightarrow$\,site\,$\rightarrow$\,design}}
\newcommand{\PSD}{\textsc{Psd}}
\newcommand{\badge}[1]{\texttt{#1}}
\newcommand{\BodyKey}{\mathrm{BodyKey}}
\newcommand{\Badge}{\mathrm{Badge}}

\setlength{\itemsep}{0.22em}
\setlength{\parsep}{0.12em}
\setlength{\topsep}{0.35em}

\title{Path, Site, Design:\\
Locating Agreement in Multi-Model Code Evaluation}
\subtitle{A comparison geometry for ranking peer patches\\
--- with formal properties, a benchmark sketch, and product patterns}
\author{Technical Proposal\\[0.35em]
{\normalsize Multi-Model Code Evaluation \& Preference Interfaces}}
\date{}

\begin{document}
\maketitle
\thispagestyle{plain}

\vspace{-1.1em}
{\color{rulegray}\hrule height 0.6pt}
\vspace{1.0em}

\begin{abstract}
Language models that write code are increasingly evaluated not only by
unit tests and reference similarity, but by \emph{human preference among
peer solutions} to the same software-engineering task. When $N$ models
each emit a multi-file patch against a shared repository base, existing
interfaces force raters either to open $N$ independent review tabs or to
treat the problem as a merge. Both miss the unit that preference
justifications actually need: \emph{the same base location under alternate
designs}.

We propose \PathSiteDesign{} (\PSD{}), a comparison geometry for
multi-model code evaluation. Paths are triaged by coverage and change
class; multi-model edits are partitioned into \emph{sites} (base-anchored
hunk groups) and \emph{designs} (body-equivalence classes at a site).
Agreement is defined only at a site, never by global token overlap. We
state the core model formally, prove elementary soundness properties
(totality of path triage, site-local agreement, stability of badges under
presentation filters), and show that a natural token-set baseline invents
false agreement that \PSD{} refuses.

We sketch \textbf{\PSD{}-Bench}, a benchmark of multi-model patch tuples
with gold site/design labels and ranking-relevant annotations, and report
an evaluation on synthetic adversaries plus a corpus of real multi-model
agent patches ($104$ paths, $136$ core sites): token-level locus conflation
is common ($17$ multi-locus paths), while true multi-model path agreement
is rare ($2$ of $104$). Finally we extract product patterns for evaluation
platforms, IDE multi-agent UIs, and preference-data pipelines---independent
of any single viewer implementation.
\end{abstract}

\vspace{0.5em}
{\color{rulegray}\hrule height 0.4pt}
\vspace{0.7em}

\tableofcontents
\clearpage

%======================================================================
\section{Introduction}
\label{sec:intro}

\subsection{From scoring code to ranking peer patches}

Benchmarks for code generation historically emphasize functional
correctness, reference match, and efficiency~\cite{zheng2024survey,effibench}.
A parallel evaluation mode is growing: \emph{preference among peer
solutions} produced by different models or agents for one task. Preference
data trains reward models and steers product ranking; written
justifications document \emph{where} and \emph{why} one design wins.

In software-engineering settings the unit of solution is often a
\textbf{patch}---a multi-file unified diff against a shared base tree---not
a single function completion. Human raters must therefore compare $N$ peer
patches. This is \emph{model code evaluation by ranking}, not integration.

\subsection{What breaks today}

Three dominant interface families fail this task in different ways:

\begin{enumerate}
  \item \textbf{Independent review tabs.} Each model gets its own ``Files
        changed'' view. Raters search by path; multi-file tasks bury the
        loci of conceptual disagreement.
  \item \textbf{Whole-solution pickers.} Best-of-$N$ product UIs show full
        transcripts or apply-one-patch buttons. Holistic preference is
        easy; site-level evidence for written justification is not.
  \item \textbf{Merge tools.} Three-way merge and multi-pane conflict
        UIs~\cite{mens2002survey,khanna2007unified} optimize
        \emph{resolution into one tree}. Ranking needs the opposite:
        preserve mutually exclusive designs at one base location.
\end{enumerate}

A fourth failure mode is algorithmic rather than interactive:
\textbf{token-set ``Shared''} over a file (or over all added lines) marks
agreement whenever models insert the same tokens, even in different
functions. That invents false consensus and poisons both UI chrome and
automatic agreement metrics.

\subsection{Contribution}

We introduce \PathSiteDesign{} (\PSD{}), a comparison geometry for
multi-model code evaluation:

\begin{enumerate}
  \item \textbf{Model.} Paths $\rightarrow$ sites $\rightarrow$ designs,
        with a total path-triage function and site relations defined from
        body keys at base anchors (Section~\ref{sec:model}).
  \item \textbf{Properties.} Proofs that triage is total, that agreement
        is site-local (hence token-set Shared is not implied), and that
        path badges remain stable under optional presentation filters
        (Section~\ref{sec:props}).
  \item \textbf{Benchmark sketch.} \PSD{}-Bench: multi-model patch tuples
        with gold labels for badges, sites, designs, and ranking-relevant
        loci (Section~\ref{sec:bench}).
  \item \textbf{Evaluation.} Synthetic adversaries and a corpus of real
        multi-model agent patches (Section~\ref{sec:eval}).
  \item \textbf{Product patterns.} Design patterns for evaluation
        platforms, multi-agent IDEs, and preference pipelines---not tied
        to one viewer (Section~\ref{sec:products}).
\end{enumerate}

\PSD{} is dual to multi-agent \emph{review} of a single change: there $N$
lenses inspect one patch; here $N$ patches face one ranking objective.

%======================================================================
\section{Related Work and Positioning}
\label{sec:related}

\begin{table}[t]
\centering
\caption{Where \PSD{} sits among nearby systems. ``Agreement unit'' is what
the tool treats as the same change when comparing alternatives.}
\label{tab:position}
\begin{tabular}{@{}p{2.6cm}p{1.6cm}p{2.0cm}p{2.0cm}p{3.4cm}@{}}
\toprule
\textbf{Family} & \textbf{Peers} & \textbf{Goal} & \textbf{Anchor} & \textbf{Agreement unit} \\
\midrule
Modern code review
  & 1 tip & accept/reject & PR diff & single proposed patch \\
Three-way merge / KDiff3
  & 2 tips & integrate & common ancestor & conflict region \\
\texttt{git range-diff}
  & 2 series & revise series & commit & patch-series interdiff \\
Structural / refactor diffs
  & 1 tip & reduce noise & AST/symbol & one tip vs base \\
Whole-solution BoN / arena UIs
  & $N$ & pick winner & none / whole & entire solution \\
Token-set ``Shared'' baselines
  & $N$ & show overlap & none & bag of added lines \\
\textbf{\PSD{} (this work)}
  & $N$ & \textbf{rank} & base hunk line & \textbf{BodyKey at site} \\
\bottomrule
\end{tabular}
\end{table}

\paragraph{Modern code review (MCR).}
MCR studies cognitive load on multi-file diffs and the gap between hunks
and design intent~\cite{bacchelli2013expectations,sadowski2018modern,wang2024roadmap}.
Almost all tooling assumes \emph{one} proposed tip. We keep the
patch-centric unit of review~\cite{janestreet2013patch} but lift it to
$N$ peers on one objective.

\paragraph{Merge and multi-version visualization.}
Classical merging constructs one result from base and
tips~\cite{mens2002survey,khanna2007unified}. Layout metaphors (side-by-side,
ancestor) transfer; the objective does not. \PSD{} never emits a merged
tree.

\paragraph{Code evaluation and preference.}
Surveys of LLM code evaluation stress execution and similarity
metrics~\cite{zheng2024survey,effibench}. Human preference interfaces, when
present, are often whole-solution or pairwise at output level. Multi-file
agent patches are rarely first-class comparison objects.

\paragraph{Gap.}
No widely used evaluation stack simultaneously (i)~takes $N$ peer patches
for one task, (ii)~anchors comparison at base hunk geometry, (iii)~defines
designs as body-equivalence classes at a site, and (iv)~refuses
token-level false agreement across unrelated locations
(Table~\ref{tab:position}).

%======================================================================
\section{Problem Formulation}
\label{sec:problem}

\begin{definition}[Multi-model patch task]
\label{def:task}
Fix a base tree $T$ and models $M=\{0,\ldots,N-1\}$. Each model $i$ emits
a set of path-indexed unified diffs $\{P_i(p)\}_{p\in\mathcal{P}_i}$ against
$T$ (including pure-add new files and pure-delete removals). A
\emph{ranking objective} asks a rater (human or automated judge) for an
order on $M$ and optional free-text justification.
\end{definition}

\begin{definition}[Comparison layer]
\label{def:layer}
A comparison layer maps $\{P_i\}$ to a navigable structure supporting
atomic judgments of the form
\[
  \text{``at path $p$, site $S$, design $D_j$ is preferable to $D_k$.''}
\]
It does \emph{not} produce a merged tree.
\end{definition}

\paragraph{Failure modes we design against.}
\begin{itemize}
  \item \textbf{Tab explosion:} $N$ full file lists; path search is manual.
  \item \textbf{False Shared:} any shared added-token multiset counts as
        agreement (e.g., both insert \texttt{return True} in different
        functions).
  \item \textbf{Merge bias:} UIs that push toward resolution rather than
        ranking.
  \item \textbf{Metric collapse:} automatic ``overlap scores'' that ignore
        base location.
\end{itemize}

%======================================================================
\section{The \PSD{} Model}
\label{sec:model}

\subsection{Coverage}

\begin{definition}[Coverage]
\label{def:coverage}
Let $P_i(p)$ be model $i$'s patch for path $p$, or empty if absent.
\[
  C(p) \;=\; \{ i \in M \mid P_i(p) \neq \emptyset \}.
\]
\end{definition}

\subsection{Hunks and body keys}

\begin{definition}[Hunk]
\label{def:hunk}
A hunk $h$ is a unified-diff region with base start $\mathrm{old}(h)\in\mathbb{N}$,
base span length $\mathrm{count}(h)\in\mathbb{N}$, and an ordered list of
typed lines (context / add / delete).
\end{definition}

\begin{definition}[Body key]
\label{def:bodykey}
\[
  \BodyKey(h)
  \;=\;
  \bigl(\,
    \tau(\ell)\;\big|\;
    \ell\text{ is add/del in }h,\;
    \mathrm{strip}_{\mathrm{CR}}(\ell)\text{ is non-blank}
  \,\bigr),
\]
where $\tau(\ell)$ is the kind-tagged text (\texttt{add:}$\ldots$ or
\texttt{del:}$\ldots$). Context lines are excluded. Blank-only add/del
noise is ignored so whitespace padding does not split designs.
\end{definition}

\subsection{Sites and designs}

\begin{definition}[Core site]
\label{def:site}
For a path $p$, parse each $P_i(p)$ into hunks. For each base line
$\ell\in\mathbb{N}$, the \emph{core site} at $\ell$ is the $N$-tuple
\[
  S_\ell
  \;=\;
  (h_0,\ldots,h_{N-1}),
  \qquad
  h_i
  \;=\;
  \begin{cases}
    \text{first hunk of $i$ with $\mathrm{old}(h)=\ell$} & \text{if any}\\
    \bot & \text{otherwise.}
  \end{cases}
\]
Sites with all slots $\bot$ are dropped. Co-location in the
\emph{normative core} is exact $\mathrm{old}(h)$ only.
\end{definition}

\begin{definition}[Designs and site relation]
\label{def:design}
Let $\mathrm{present}(S_\ell)=\{i\mid h_i\neq\bot\}$. Partition
$\mathrm{present}(S_\ell)$ by equality of $\BodyKey(h_i)$. Each block is a
\emph{design}. The site relation $\rho(S_\ell)$ is:
\begin{center}
\begin{tabular}{@{}ll@{}}
\toprule
$\rho$ & condition \\
\midrule
\texttt{unique} & $|\mathrm{present}|=1$ \\
\texttt{same} & one design, $\ge 2$ models \\
\texttt{pair} & $N=3$ specialization: sizes $(2,1)$ \\
\texttt{diverge} & $\ge 2$ designs (and not the pair pattern) \\
\bottomrule
\end{tabular}
\end{center}
For general $N$, the relation is the multiset of design sizes (a partition
of $|\mathrm{present}|$); \texttt{pair} is the product name for type
$2{+}1$ when $N=3$.
\end{definition}

\begin{remark}[Why exact $\mathrm{old}$ is normative]
Peer patches against the \emph{same} base do not shift each other's
\texttt{old\_start}. Misalignment comes from different hunkification of
one conceptual region, not tip drift. Optional presentation layers may
coalesce overlapping spans when BodyKeys do not conflict; they must not
redefine the core (Section~\ref{sec:props}).
\end{remark}

\subsection{Path triage}

\begin{definition}[Path badge]
\label{def:badge}
$\Badge(p)$ is defined by the total decision order:
\begin{enumerate}
  \item If $C(p)=\emptyset$: omit $p$.
  \item Else if every present patch is a pure full-file delete: $\badge{del}$.
  \item Else if every present patch is a pure full-file add: $\badge{new}$.
  \item Else if $|C(p)|=1$: $\badge{solo}$ (single-model \emph{edit}).
  \item Else if all present patches are equal after normalization
        (strip CR; drop ``No newline at end of file'' markers): $\badge{same}$.
  \item Else let $\mathcal{S}$ be the multi-present core sites
        ($|\mathrm{present}|\ge 2$). If any has $\rho\in\{\texttt{diverge},\texttt{pair}\}$:
        $\badge{diff}$. Else if any has $\rho=\texttt{same}$: $\badge{share}$.
        Else: $\badge{diff}$.
\end{enumerate}
Normalization is encoding hygiene, not semantic equivalence.
\end{definition}

Full-file $\badge{new}$/$\badge{del}$ take precedence over $\badge{same}$:
identical multi-model pure-adds remain $\badge{new}$; pure deletes remain
$\badge{del}$. Path $\badge{share}$ requires multi-present agreement
\emph{without} multi-present diverge/pair---partial agreement does not
upgrade the path.

\begin{example}[False Shared refused]
\label{ex:false-shared}
Models $A,B,C$ all insert the token \texttt{return True}, but $A$ and $C$
at function $f_1$ (line $10$) while $B$ inserts it only at $f_2$ (line $80$),
and bodies otherwise differ. Token-set intersection is nonempty, so a
bag-of-lines baseline reports ``Shared.'' Core sites are $S_{10}$ and
$S_{80}$ with relations \texttt{pair}/\texttt{diverge} as bodies dictate;
no path-level Shared is invented. Empirically this adversary yields
$\Badge=\badge{diff}$ with two pair sites (Section~\ref{sec:eval}).
\end{example}

\subsection{Presentation (non-normative)}

Optional layers may (i)~drop non-signal sites (comment/punct filters),
(ii)~coalesce span-overlapping sites without BodyKey conflict,
(iii)~collapse identical fingerprints, (iv)~present multi-design sites as
tabs under one linear anchor, (v)~compress near-identical pure-add full
files for display. These affect \emph{detail UI density only}. They never
recompute $\Badge(p)$ (Policy~A).

%======================================================================
\section{Properties}
\label{sec:props}

We treat patches as immutable inputs and $\BodyKey$ as a pure function of
hunk lines as in Definition~\ref{def:bodykey}.

\begin{proposition}[Totality of triage]
\label{prop:total}
For every path $p$ with $C(p)\neq\emptyset$, exactly one badge in
$\{\badge{del},\badge{new},\badge{solo},\badge{same},\badge{share},\badge{diff}\}$
is assigned.
\end{proposition}

\begin{proof}
The decision order in Definition~\ref{def:badge} is a sequence of mutually
exclusive guards covering all nonempty $C(p)$. Pure-delete and pure-add are
checked before solo/identical; identical multi-edits before site analysis;
the final site branch's outcomes all map to $\badge{share}$ or $\badge{diff}$.
No fall-through remains.
\end{proof}

\begin{proposition}[Site-local agreement]
\label{prop:local}
Suppose models $i$ and $j$ both contain an add-line with identical text
$t$ in path $p$, but only in hunks $h_i,h_j$ with
$\mathrm{old}(h_i)\neq\mathrm{old}(h_j)$. Then there is no core site $S$
at which $i$ and $j$ share a design solely because of $t$ appearing in both
patches. In particular, a non-empty global intersection of added-token sets
does not imply any site with $\rho=\texttt{same}$.
\end{proposition}

\begin{proof}
Core sites are keyed only by $\mathrm{old}(\cdot)$
(Definition~\ref{def:site}). Tokens at distinct base starts lie in distinct
sites $S_{\ell}$ and $S_{\ell'}$. Within one site, designs depend on the
full $\BodyKey$ of each present hunk (Definition~\ref{def:design}), not on
membership of $t$ in a path-level set. Therefore path-level token
intersection cannot create a $\texttt{same}$ site, and by
Definition~\ref{def:badge} cannot by itself produce $\badge{share}$ or
$\badge{same}$.
\end{proof}

\begin{corollary}[Token-set baseline is strictly coarser]
\label{cor:token}
Let $\mathrm{TokShared}(p)$ hold when
$\bigcap_{i\in C(p)}\mathrm{Adds}(P_i(p))\neq\emptyset$ with $|C(p)|\ge 2$.
There exist patch tuples for which $\mathrm{TokShared}(p)$ is true while
$\Badge(p)\notin\{\badge{share},\badge{same}\}$ and no core site has
$\rho=\texttt{same}$. (Example~\ref{ex:false-shared}.)
\end{corollary}

\begin{proposition}[Policy~A: badge stability under presentation]
\label{prop:policyA}
Let $\mathcal{S}_{\mathrm{core}}(p)$ be the core site list and
$F$ any presentation map (filter/coalesce) such that:
\begin{enumerate}
  \item $F$ only drops sites or merges sites whose filled slots have
        equal $\BodyKey$ pairwise per model;
  \item $F$ never invents a new filled slot with a BodyKey not present in
        some core site for that model.
\end{enumerate}
If $\Badge(p)$ is computed from $\mathcal{S}_{\mathrm{core}}$ as in
Definition~\ref{def:badge}, then replacing $\mathcal{S}_{\mathrm{core}}$ by
$F(\mathcal{S}_{\mathrm{core}})$ inside the badge function can change the
badge only if an implementation \emph{violates} Policy~A. Under Policy~A,
$\Badge(p)$ is independent of $F$.
\end{proposition}

\begin{proof}
Immediate from the definition of Policy~A: the badge function's site
branch is evaluated on $\mathcal{S}_{\mathrm{core}}$ only. The constraints
on $F$ describe the intended presentation layer (no invented bodies); they
are not required for independence once the badge is defined to ignore $F$.
\end{proof}

\begin{remark}
Proposition~\ref{prop:policyA} is an interface contract as much as a
theorem: products that recompute badges after filtering can flip
$\badge{share}\leftrightarrow\badge{diff}$ when a lone agreeing site is
dropped. That variant needs its own tests; it is not \PSD{}-core.
\end{remark}

\begin{lemma}[Design partition]
\label{lem:partition}
At each site $S$, the designs form a partition of $\mathrm{present}(S)$.
Hence every present model appears in exactly one design card emission.
\end{lemma}

\begin{proof}
Designs are equivalence classes under equality of $\BodyKey(h_i)$ on the
finite set $\mathrm{present}(S)$.
\end{proof}

\begin{proposition}[Fallback completeness]
\label{prop:fallback}
If a multi-model edit path yields no core sites usable for detail
emission, the comparison layer still exposes $|C(p)|$ per-model full-patch
designs. Ranking remains possible; only site alignment is lost.
\end{proposition}

\begin{proof}
By construction of the fallback (normative): emit one card per present
model with patch text $P_i(p)$. Lemma~\ref{lem:partition} is vacuous;
coverage $C(p)$ is unchanged.
\end{proof}

\paragraph{What we do not claim.}
$\BodyKey$ equality is textual, not semantic or behavioral. Exact
$\mathrm{old}$ co-location is incomplete w.r.t.\ human ``same function''
judgments when diff tools split regions differently. \PSD{} makes that
incompleteness explicit rather than papering over it with token sets.

%======================================================================
\section{Algorithms (Core)}
\label{sec:algo}

We give the normative core compactly. Indices are $0$-based; product UIs
may map them to letters.

\begin{lstlisting}[caption={Normative core: sites, body keys, path badge.},label={lst:core}]
function BodyKey(h):
    return tuple( kind:stripCR(text)
                  for kind,text in add/del lines of h
                  if strip(text) nonempty )

function BuildSites(patches):          # exact old_start only
    by_start := map ell -> array[N] of (hunk|bot)
    for i in models:
        for h in ParseHunks(patches[i]):
            if by_start[h.old_start][i] is bot:
                by_start[h.old_start][i] := h
    return [ Site(ell, slots, Classify(slots))
             for ell,slots in by_start sorted ]

function Classify(slots):
    present := {i | slots[i] != bot}
    if |present| <= 1: return unique
    keys := { BodyKey(slots[i]) | i in present }
    if |keys| == 1: return same
    # general N: return partition-type(keys);
    # N=3 product: pair if type 2+1 else diverge
    ...

function Badge(patches):               # Policy A: core sites only
    present := {i | patch i present}
    if all pure-delete: return del
    if all pure-add:    return new
    if |present| == 1:  return solo
    if all Normalize(patch i) equal: return same
    multi := { s in BuildSites(patches) | |present(s)| >= 2 }
    if any multi has relation in {diverge, pair}: return diff
    if any multi has relation same:               return share
    return diff
\end{lstlisting}

\paragraph{List geometry (evaluation UX).}
Sort paths by kind rank
($\badge{diff}\prec\badge{share}\prec\badge{new}\prec\badge{del}\prec\badge{same}\prec\badge{solo}$),
then coverage rank, then path. Group consecutive rows that share
$(\Badge, C(p))$ under one coverage header so ``who touched this'' is
stated once per block. Detail views show one linear site sequence; when a
site has $\ge 2$ designs, tabs switch bodies without stacking $N$ full
walls.

%======================================================================
\section{\PSD{}-Bench: A Benchmark Sketch}
\label{sec:bench}

Evaluation of multi-model code systems needs more than pass@k. We propose
\textbf{\PSD{}-Bench} as a dataset-and-metrics layer that any evaluation
platform can adopt.

\subsection{Instance schema}

Each instance is a tuple
\[
  \bigl(T,\; q,\; (P_0,\ldots,P_{N-1}),\; G,\; R^\star\bigr)
\]
where $T$ is a base revision (or minimal virtual tree), $q$ a natural-language
task, $P_i$ peer patches, $G$ gold comparison labels, and $R^\star$ an optional
gold ranking / preference with justification spans.

Gold $G$ includes, per path:
\begin{itemize}
  \item $\Badge^\star(p)$ and coverage $C(p)$;
  \item core sites $\{(\ell,\rho^\star,\text{design partition}^\star)\}$;
  \item flags: \texttt{false\_token\_shared} (token baseline errs),
        \texttt{hunk\_misalign} (human same-region, distinct $\mathrm{old}$),
        \texttt{fallback}.
\end{itemize}

\subsection{Splits}

\begin{enumerate}
  \item \textbf{Synthetic control.} Hand-crafted adversaries (false Shared;
        identical pure-new vs $\badge{same}$; pair vs diverge; CRLF; partial
        agreement $\Rightarrow$ $\badge{diff}$).
  \item \textbf{Agent corpus.} Real multi-model patches on shared bases
        (anonymized), gold labels from expert annotators using the \PSD{}
        codebook.
  \item \textbf{Stress.} Large new-file galleries, renames, binaries,
        $N\in\{2,3,5\}$.
\end{enumerate}

\subsection{Metrics}

\begin{table}[t]
\centering
\caption{Proposed \PSD{}-Bench metrics.}
\label{tab:metrics}
\begin{tabular}{@{}lp{9.2cm}@{}}
\toprule
\textbf{Metric} & \textbf{Definition} \\
\midrule
Badge macro-F1 & Per-class F1 of $\Badge$ vs $\Badge^\star$ \\
Site anchor recall & Fraction of gold sites matched within exact
  $\mathrm{old}$ (core) or span-IoU $\ge\tau$ (presentation) \\
Design ARI & Adjusted Rand Index of design partitions at matched sites \\
False Shared rate & Fraction of paths where token baseline claims Shared
  but gold has no multi-present $\texttt{same}$ site \\
Fallback rate & Multi-model edit paths with empty core sites \\
Rank locality & Among pairs of models, fraction of gold preference
  justifications spannable by a single site or path (upper bound on
  UI sufficiency) \\
\bottomrule
\end{tabular}
\end{table}

\subsection{Baselines}

\begin{itemize}
  \item Independent file lists (no alignment).
  \item Token-set Shared / Jaccard on added lines.
  \item Pairwise \texttt{diff} of tips without base sites.
  \item \PSD{}-core; \PSD{}-core$+$span coalesce (presentation).
\end{itemize}

\subsection{Why this is unique as an evaluation artifact}

Existing code benchmarks score \emph{solutions}. Preference arenas score
\emph{whole outputs}. Merge suites score \emph{integration}.
\PSD{}-Bench scores the \textbf{comparison geometry} itself---the missing
middle layer between raw patches and a ranking label. That layer is what
product UIs and automated judges both need if justifications are to be
faithful.

%======================================================================
\section{Evaluation}
\label{sec:eval}

We evaluate the \emph{core geometry} (not a particular pixel UI) on
(1)~synthetic adversaries that encode the design principles, and
(2)~a convenience corpus of real multi-model agent patch dumps available
to this study.

\subsection{Synthetic conformance}

Table~\ref{tab:conform} summarizes control cases. All are decidable from
Definitions~\ref{def:coverage}--\ref{def:badge}; an open-source conformance
suite can lock them as fixtures.

\begin{table}[t]
\centering
\caption{Core conformance scenarios (synthetic).}
\label{tab:conform}
\begin{tabular}{@{}clp{6.2cm}@{}}
\toprule
\# & Scenario & Expected core outcome \\
\midrule
1 & Multi-model identical edit & $\badge{same}$, one design \\
2 & Multi-model identical pure-new & $\badge{new}$ (not $\badge{same}$) \\
3 & Solo pure-delete & $\badge{del}$, removed body \\
4 & Bodies $(2{+}1)$ at one $\mathrm{old}$ & site \texttt{pair}; path $\badge{diff}$ \\
5 & One same site $+$ one diverge site & path $\badge{diff}$ \\
6 & Same token at two functions & two sites; no Shared \\
7 & No aligned sites & fallback cards; $\badge{diff}$ \\
8 & Mid-file pure-adds & BodyKey at each $\mathrm{old}$ only \\
9 & CRLF vs LF identical edits & $\badge{same}$ via Normalize \\
10 & All sites unique, coverage ABC & $\badge{diff}$, ABC header \\
11 & Same body, two $\mathrm{old}$ values & two sites (no core coalesce) \\
12 & $N{=}2$, two bodies & site \texttt{diverge} \\
\bottomrule
\end{tabular}
\end{table}

\paragraph{Adversary result (Example~\ref{ex:false-shared}).}
On the cross-placed \texttt{return True} triple, token intersection is
$\{\texttt{return True}\}$ (baseline would claim Shared) while \PSD{}
reports $\badge{diff}$ with sites at lines $10$ and $80$, each
\texttt{pair}---Proposition~\ref{prop:local} in action.

\subsection{Corpus study}

\paragraph{Material.}
Six multi-model patch tasks ($N=3$ each) drawn from agent ranking dumps
and public fixtures used in tooling tests: $104$ union paths after
comparison, $28$ multi-model paths, $136$ core sites.

\begin{table}[t]
\centering
\caption{Path badge distribution on the multi-model corpus ($104$ paths).}
\label{tab:corpus-badges}
\begin{tabular}{@{}lrr@{}}
\toprule
\textbf{Badge} & \textbf{Count} & \textbf{Share} \\
\midrule
\badge{solo}  & 64 & 61.5\% \\
\badge{diff}  & 22 & 21.2\% \\
\badge{new}   & 16 & 15.4\% \\
\badge{share} & 1  & 1.0\% \\
\badge{same}  & 1  & 1.0\% \\
\badge{del}   & 0  & 0.0\% \\
\midrule
Total & 104 & 100\% \\
\bottomrule
\end{tabular}
\end{table}

\begin{table}[t]
\centering
\caption{Site relations on the same corpus ($136$ core sites).}
\label{tab:corpus-sites}
\begin{tabular}{@{}lrr@{}}
\toprule
\textbf{Relation} & \textbf{Count} & \textbf{Share} \\
\midrule
\texttt{unique} (\texttt{only}) & 91 & 66.9\% \\
\texttt{diverge} & 35 & 25.7\% \\
\texttt{pair} & 6 & 4.4\% \\
\texttt{same} & 4 & 2.9\% \\
\bottomrule
\end{tabular}
\end{table}

\paragraph{Findings.}

\begin{enumerate}
  \item \textbf{True path-level multi agreement is rare.} Only $2/104$
        paths are $\badge{share}$ or $\badge{same}$. Most multi-model
        contact is partial or conflicting---exactly where site-level
        designs matter.
  \item \textbf{Token locus conflation is common.} On $17$ paths the same
        nonempty add-token appears at $\ge 2$ distinct core starts under
        multi-model coverage. A file-level token Shared baseline would
        conflate loci on $\approx 16\%$ of all paths and a large fraction
        of multi-model paths---Corollary~\ref{cor:token} at corpus scale.
  \item \textbf{Solo dominates raw path count.} $61.5\%$ of paths are
        single-model. List geometry should not spend equal chrome on every
        row; coverage grouping and kind order push $\badge{diff}$ first.
  \item \textbf{Site mass is real.} $136$ sites on $104$ paths
        ($\approx 1.3$ sites/path overall; much higher on heavy
        multi-model tasks) justify site navigation over whole-file only.
\end{enumerate}

\paragraph{Threats to validity.}
Convenience corpus (not a published benchmark); $N=3$ only; no inter-rater
study of gold labels; agent dumps may not match production traffic. These
motivate \PSD{}-Bench rather than replace it.

%======================================================================
\section{Product Patterns for Future Systems}
\label{sec:products}

\PSD{} is a comparison geometry, not a single application. We extract
patterns for builders of evaluation platforms, multi-agent IDEs, and data
pipelines.

\subsection{Pattern P1 --- Comparison object, not only score}

\textbf{Problem.} Leaderboards store a scalar preference; justifications
lose structure.\\
\textbf{Design.} Persist $(p,S,D)$ references alongside rankings. Reward
models and audit tools can attend to site-scoped evidence.\\
\textbf{Why unique.} Elevates the comparison layer to a first-class
evaluation artifact (Section~\ref{sec:bench}).

\subsection{Pattern P2 --- Rank surface, merge surface}

\textbf{Problem.} Teams reuse merge widgets for peer agents.\\
\textbf{Design.} Separate modes: \emph{Rank} preserves exclusive designs
(tabs/cards); \emph{Merge} is opt-in and never the default for preference
collection.\\
\textbf{Why unique.} Makes the objective mismatch explicit in UX copy and
API names.

\subsection{Pattern P3 --- Progressive disclosure}

\textbf{Problem.} Flat $N$-pane file views do not scale.\\
\textbf{Design.} path triage $\rightarrow$ coverage groups $\rightarrow$
site anchors $\rightarrow$ design tabs. Keyboard jumps to each model's full
patch remain available.\\
\textbf{Why unique.} Aligns navigation cost with disagreement mass
(corpus: few $\badge{share}$, many $\badge{diff}$/solo).

\subsection{Pattern P4 --- Policy~A metrics}

\textbf{Problem.} Presentation filters quietly change agreement dashboards.\\
\textbf{Design.} Report badges and agreement rates from core sites only;
publish filter configs separately.\\
\textbf{Why unique.} Prevents ``improving'' agreement by dropping diverge
sites.

\subsection{Pattern P5 --- Anti-metrics}

\textbf{Problem.} Jaccard-on-adds looks like quality of consensus.\\
\textbf{Design.} Ban path-level token Shared as a primary metric; if shown,
label as ``token overlap (not site agreement)'' and pair with False Shared
rate from \PSD{}-Bench.\\
\textbf{Why unique.} Stops a seductive but invalid KPI.

\subsection{Pattern P6 --- Judge models with site prompts}

\textbf{Problem.} LLM-as-judge sees concatenated patches and hallucinates
shared intent.\\
\textbf{Design.} Prompt judges per site with the design bodies only; ask
for local preference; aggregate with explicit rules.\\
\textbf{Why unique.} Transfers \PSD{} from human UI to automated evaluation.

\subsection{Pattern P7 --- N-ready slots, named algebras}

\textbf{Problem.} Hard-coded ABC pair logic blocks $N{>}3$.\\
\textbf{Design.} Store sites as length-$N$ slot arrays; expose relation as
partition type; keep ``pair'' as sugar for $2{+}1$ when $N=3$.\\
\textbf{Why unique.} Separates geometry from product vocabulary.

\subsection{Reference mapping}

A keyboard-first shared-file compare view in ranking tools (e.g.\ open-source
Starfleet Control) is one \emph{instance} of P2--P5. The geometry does not
depend on that instance; any evaluation console or IDE multi-agent panel
can implement the same contracts.

%======================================================================
\section{Discussion}
\label{sec:discussion}

\subsection{Implications for model code evaluation}

If the field cares about preference data on agentic coding, it must care
about \textbf{where} models diverge. Pass@k does not say; whole-output
arenas barely say; \PSD{} makes the locus explicit. That benefits human
raters and structured automatic judges alike.

\subsection{Relation to structural diff}

AST-aware anchors could improve co-location when line numbers drift.
They are compatible as an optional co-location oracle \emph{above} the
core, subject to Policy~A. We keep line-level core first because unified
diffs are the lingua franca of agent tools today.

\subsection{Limitations}

\begin{itemize}
  \item Textual BodyKeys miss semantic equivalence (rename-only, formatting
        tools, import reorder).
  \item Exact $\mathrm{old}$ under-merges some human ``same edit'' pairs.
  \item Corpus evaluation is exploratory; \PSD{}-Bench is a proposal.
  \item General-$N$ relation UX is underspecified beyond partition types.
\end{itemize}

\subsection{Rejected alternatives}

\begin{itemize}
  \item Token-set Shared over files or mid-file pure-add sites.
  \item Pairwise-only screens without simultaneous multi-design sites.
  \item Always-on $N$-pane full files as the primary surface.
  \item Merge layout as the default preference surface.
  \item Path $\badge{share}$ on any single agreeing site (over-promotes
        partial consensus).
  \item Identical multi-new as $\badge{same}$ (wrong mental category).
\end{itemize}

%======================================================================
\section{Conclusion}
\label{sec:conclusion}

Multi-model code evaluation needs a comparison geometry between raw
patches and a ranking label. \PathSiteDesign{} defines that geometry:
coverage-aware paths, base-anchored sites, and body-equivalence designs,
with agreement that is site-local by construction. We proved elementary
soundness properties, sketched \PSD{}-Bench, and showed on synthetic and
real multi-model patches that false token agreement is common while true
path-level multi agreement is rare---so site-level designs are not a
luxury.

Future work: publish \PSD{}-Bench with multi-annotator gold; generalize
relation algebras for $N{>}3$; study LLM-as-judge accuracy with site-scoped
prompts (Pattern~P6); and measure rater time-to-justification against
independent-tab baselines.

%======================================================================
\section*{Notation}

\begin{center}
\begin{tabular}{@{}ll@{}}
\toprule
\textbf{Symbol} & \textbf{Meaning} \\
\midrule
$M$ & model indices \\
$P_i(p)$ & unified patch of model $i$ for path $p$ \\
$C(p)$ & coverage set \\
$\BodyKey(h)$ & ordered non-blank add/del lines of hunk $h$ \\
$S_\ell$ & core site at base line $\ell$ \\
$\rho(S)$ & site relation (partition type / product name) \\
$\Badge(p)$ & path triage class \\
\PSD{} & path\,$\rightarrow$\,site\,$\rightarrow$\,design geometry \\
\bottomrule
\end{tabular}
\end{center}

%======================================================================
\begin{thebibliography}{99}

\bibitem{bacchelli2013expectations}
A.~Bacchelli and C.~Bird.
\newblock Expectations, outcomes, and challenges of modern code review.
\newblock In \emph{ICSE}, 2013.

\bibitem{effibench}
Y.~Huang et~al.
\newblock EffiBench-X: A multi-language benchmark for measuring efficiency of
LLM-generated code.
\newblock arXiv:2505.13004, 2025.

\bibitem{janestreet2013patch}
Jane Street.
\newblock Patch review vs.\ diff review, revisited.
\newblock Engineering blog, 2013.
\newblock \url{https://blog.janestreet.com/patch-review-vs-diff-review-revisited/}

\bibitem{khanna2007unified}
S.~Khanna, K.~Kunal, and B.~C.~Pierce.
\newblock A formal investigation of diff3.
\newblock In \emph{FSTTCS}, 2007.

\bibitem{mens2002survey}
T.~Mens.
\newblock A state-of-the-art survey on software merging.
\newblock \emph{IEEE TSE}, 28(5), 2002.

\bibitem{sadowski2018modern}
C.~Sadowski, E.~S{\"o}derberg, L.~Church, M.~Sipko, and A.~Bacchelli.
\newblock Modern code review: A case study at Google.
\newblock In \emph{ICSE-SEIP}, 2018.

\bibitem{wang2024roadmap}
X.~Wang et~al.
\newblock A roadmap for modern code review: Challenges and opportunities.
\newblock arXiv:2405.18216, 2024.

\bibitem{zheng2024survey}
Z.~Zheng et~al.
\newblock A survey on evaluating large language models in code generation tasks.
\newblock arXiv:2408.16498, 2024.

\end{thebibliography}

\end{document}
